% NeurIPS-style method draft for the OmniSearch UAV survivor-search reward system.
% This file is intentionally standalone so it can be compiled directly, or its
% sections can be copied into a larger NeurIPS manuscript.
\documentclass{article}

\usepackage[margin=1in]{geometry}
\usepackage{amsmath}
\usepackage{amssymb}
\usepackage{booktabs}
\usepackage{enumitem}
\usepackage{hyperref}
\usepackage{xcolor}

\title{Confidence-Aware Multi-UAV Reinforcement Learning for Survivor Search}
\author{OmniSearch UAV Diagnostics}
\date{}

\newcommand{\uav}{\mathrm{uav}}
\newcommand{\conf}{C}
\newcommand{\unc}{U}
\newcommand{\grid}{\Omega}
\newcommand{\foot}{F}
\newcommand{\pos}{\mathbf{p}}
\newcommand{\act}{\mathbf{a}}
\newcommand{\move}{\Delta \mathbf{p}}
\newcommand{\vect}{\mathbf{v}}
\newcommand{\eps}{\varepsilon}

\begin{document}
\maketitle

\begin{abstract}
We present a confidence-aware reward and observation design for multi-UAV
reinforcement learning in probabilistic survivor search. The central observation
from our diagnostics is that binary area coverage is not equivalent to survivor
search quality: agents can achieve high geometric coverage while survivors are
only weakly exposed near the edge of the camera footprint, under adverse cover,
or during short accidental flyovers. We therefore replace binary ``visited''
coverage as the primary learning signal with an inspection-confidence field that
tracks the cumulative probability that a survivor at each map cell would have
been detected. The resulting reward gives dense credit for improving detection
confidence, permits useful revisits of low-confidence areas, suppresses wasteful
revisits of already well-inspected areas, and assigns marginal credit in
multi-UAV settings. We combine this objective with uncertainty-frontier
observations, confidence-gated movement rewards, survivor scouting bonuses, and
soft boundary penalties. The design is intended for the communication-available
UAV diagnostic setting with three UAVs, five survivors, no UGVs, and a shared
coverage memory.
\end{abstract}

\section{Problem Setting}

We consider cooperative search over a bounded two-dimensional terrain
\(\grid \subset \mathbb{R}^2\). At each time step \(t\), \(N\) UAVs choose
continuous planar actions and move according to the VMAS dynamics. Survivors are
stationary and initially unknown. A UAV can scout a survivor only if the survivor
lies inside the UAV camera footprint, but detection is probabilistic and depends
on footprint-relative position, terrain/cover class, smoke, and flight quality.
The goal is to maximize survivor recall within a fixed horizon while maintaining
efficient search coverage and avoiding boundary-stuck or redundant trajectories.

The environment returns an individual scalar reward \(r_{i,t}\) for each UAV
\(i\). HAPPO optimizes these per-agent rewards with a centralized critic and
decentralized actors. Team terms are implemented by adding the same scalar
quantity to each UAV reward, while individual terms remain attached to the UAV
that produced the corresponding behavior. This preserves the current successful
inductive bias of individual coverage/search rewards while allowing explicit
mission-level credit assignment.

\section{Diagnostic Motivation}

The previous UAV diagnostic policy achieved high mean geometric coverage and
competitive recall compared with lawnmower search, but diagnostics revealed
three persistent failure modes. First, final binary coverage overestimated
search quality: missed survivors were often technically covered but only at the
edge of the footprint or with low cumulative detection probability. Second,
historical overlap penalties were hard to tune because some overlap is physically
unavoidable when the footprint radius is large relative to the UAV step length.
Third, frontier alignment to binary uncovered regions became less informative
late in an episode, when remaining useful regions were sparse and revisiting
low-confidence cells could still be valuable.

These observations imply that the correct state variable is not whether a cell
has ever been covered, but how confidently that cell has been inspected. The
final system therefore treats search as probabilistic belief improvement. Binary
coverage remains a reporting metric and visualization aid, but it is no longer
the main objective used to shape behavior.

\section{Confidence Map}

\subsection{Instantaneous Detection Probability}

Let \(x \in \grid\) be a grid cell and let \(\foot_{i,t}\) denote the camera
footprint of UAV \(i\) at time \(t\). The instantaneous probability that UAV
\(i\) would detect a survivor at cell \(x\) is
\begin{equation}
    p_{i,t}(x)
    =
    \mathbb{I}[x \in \foot_{i,t}]
    \; q_{\mathrm{range}}(d_{i,t}(x))
    \; q_{\mathrm{cover}}(x)
    \; q_{\mathrm{smoke}}(x,t)
    \; q_{\mathrm{flight}}(i,t),
    \label{eq:instant_detection}
\end{equation}
where \(d_{i,t}(x) \in [0,1]\) is the normalized distance from the footprint
center to cell \(x\). We use the quadratic range model
\begin{equation}
    q_{\mathrm{range}}(d)
    =
    1 - (1 - q_{\min}) d^2,
    \qquad q_{\min}=0.40,
    \label{eq:range_model}
\end{equation}
so central passes retain high detection probability and edge passes remain
useful but substantially weaker. The cover factors are RGB-only detection
factors derived from the terrain cover class. In the current realistic RGB
setting they are
\begin{equation}
    q_{\mathrm{cover}} \in
    (1.00, 0.95, 0.75, 0.55, 0.45, 0.90),
    \label{eq:cover_factors}
\end{equation}
for the simulator's ordered cover categories. Importantly, Eq.~\ref{eq:instant_detection}
does not depend on true survivor positions and is therefore not a privileged
reward.

\subsection{Cumulative Inspection Confidence}

We maintain a shared confidence map \(\conf_t(x) \in [0,1]\), interpreted as the
probability that a survivor located at \(x\) would have been detected by at
least one UAV before time \(t\). The corresponding residual uncertainty is
\begin{equation}
    \unc_t(x) = 1 - \conf_t(x).
\end{equation}
Given all UAV footprints at step \(t\), the team confidence update is
\begin{equation}
    \conf_{t+1}(x)
    =
    1
    -
    (1-\conf_t(x))
    \prod_{i=1}^{N} (1 - p_{i,t}(x)).
    \label{eq:confidence_update}
\end{equation}
This update naturally handles repeated passes: revisiting a poorly inspected
cell can still increase confidence, while repeated passes over a high-confidence
cell produce little or no gain.

\section{Reward Design}

The UAV reward is the sum of individual confidence-improvement terms, small
guidance terms, mission rewards, and boundary penalties:
\begin{equation}
\begin{aligned}
    r_{i,t}
    =
    &\; r^{\mathrm{conf}}_{i,t}
    + \lambda_{\mathrm{team}} r^{\mathrm{team\text{-}conf}}_{t}
    + r^{\mathrm{move}}_{i,t}
    + r^{\mathrm{frontier}}_{i,t} \\
    &+ r^{\mathrm{scout}}_{i,t}
    + r^{\mathrm{team\text{-}scout}}_{t}
    + r^{\mathrm{all\text{-}found}}_{t}
    - c^{\mathrm{outside}}_{i,t}
    - c^{\mathrm{outward}}_{i,t}.
    \label{eq:total_reward}
\end{aligned}
\end{equation}
All dense terms are normalized to be comparable per UAV per step. The scout and
mission-completion terms are sparse and are intentionally kept moderate so they
define the task objective without overwhelming the search-shaping signal.

\subsection{Marginal Confidence Gain}

For each UAV, we compute the counterfactual marginal contribution of that UAV's
footprint to the team confidence update. Let
\begin{equation}
    \conf^{\mathrm{all}}_{t+1}(x)
    =
    1
    -
    (1-\conf_t(x))
    \prod_{j=1}^{N} (1 - p_{j,t}(x))
\end{equation}
and let
\begin{equation}
    \conf^{-i}_{t+1}(x)
    =
    1
    -
    (1-\conf_t(x))
    \prod_{j \ne i} (1 - p_{j,t}(x)).
\end{equation}
The marginal confidence gain assigned to UAV \(i\) is
\begin{equation}
    \Delta \conf_{i,t}(x)
    =
    \conf^{\mathrm{all}}_{t+1}(x)
    -
    \conf^{-i}_{t+1}(x).
    \label{eq:marginal_gain}
\end{equation}
This counterfactual credit assignment solves two problems simultaneously. It
avoids double counting when two UAVs inspect the same cell in the same step, and
it avoids harsh hand-tuned inter-UAV overlap penalties for overlap that was not
actually redundant.

\subsection{Confidence-Weighted Coverage Reward}

The primary dense reward is the confidence-weighted marginal gain
\begin{equation}
    r^{\mathrm{conf}}_{i,t}
    =
    k_{\mathrm{conf}}
    \frac{1}{|\grid|}
    \sum_{x \in \grid}
    w(\conf_t(x)) \Delta \conf_{i,t}(x),
    \label{eq:confidence_reward}
\end{equation}
where
\begin{equation}
    w(\conf) = \eps + (1-\conf)^\gamma.
    \label{eq:confidence_weight}
\end{equation}
The weight \(w\) gives the highest value to uninspected cells, moderate value to
low-confidence revisits, and negligible value to repeatedly inspected cells. In
practice we use \(\eps > 0\) to avoid making late-episode cleanup entirely
rewardless, and \(\gamma > 1\) to focus the policy on genuinely uncertain
regions.

The team confidence reward is
\begin{equation}
    r^{\mathrm{team\text{-}conf}}_t
    =
    \frac{1}{N}
    \sum_{i=1}^{N} r^{\mathrm{conf}}_{i,t},
    \label{eq:team_confidence_reward}
\end{equation}
which is added to every UAV reward with weight \(\lambda_{\mathrm{team}}\). This
encourages cooperation without removing individual accountability for useful
inspection.

\subsection{Confidence-Gated Movement Reward}

The movement reward encourages fast useful sweeping rather than speed for its
own sake. Let \(d_{\max}\) be the maximum physical displacement in one step and
let
\begin{equation}
    s_{i,t} =
    \min \left( \frac{\|\move_{i,t}\|_2}{d_{\max}}, 1 \right)
\end{equation}
be normalized progress. We define a normalized confidence gain
\begin{equation}
    g_{i,t}
    =
    \frac{
        \sum_{x \in \grid} w(\conf_t(x)) \Delta \conf_{i,t}(x)
    }{
        \sum_{x \in \grid} w(\conf_t(x)) p^{\max}_{i,t}(x) + \eta
    },
    \label{eq:normalized_conf_gain}
\end{equation}
where \(p^{\max}_{i,t}\) is the best confidence gain available to UAV \(i\) under
a one-step max-speed candidate move over a fixed action stencil, and \(\eta\) is
a small constant. The movement reward is
\begin{equation}
    r^{\mathrm{move}}_{i,t}
    =
    k_{\mathrm{move}} s_{i,t} \mathrm{clip}(g_{i,t}, 0, 1).
    \label{eq:move_reward}
\end{equation}
This term rewards large steps only when they generate inspection value. It does
not pay for spinning, boundary pushing, or repeated motion over high-confidence
areas.

\subsection{Uncertainty Frontier Reward}

The frontier signal is computed from the uncertainty map \(\unc_t = 1-\conf_t\)
rather than from binary uncovered cells. Around each UAV we divide a local disk
into \(S=8\) angular sectors and compute a sector score
\begin{equation}
    \sigma_{i,t}^{(s)}
    =
    \frac{
        \sum_{x \in \mathcal{S}_{i,t}^{(s)}} w(\conf_t(x)) \unc_t(x)
    }{
        |\mathcal{S}_{i,t}^{(s)}| + \eta
    }.
    \label{eq:sector_score}
\end{equation}
The top \(K=2\) sectors are exposed to the policy as frontier vectors
\(\vect_{i,t}^{(1)}, \vect_{i,t}^{(2)}\), each containing a unit direction,
normalized distance to the sector's uncertainty mass, and the normalized sector
score. Providing two directions avoids the weak compromise produced by a single
centroid direction and lets the policy choose between multiple useful search
regions.

The frontier reward is a small alignment bonus:
\begin{equation}
    r^{\mathrm{frontier}}_{i,t}
    =
    k_{\mathrm{frontier}}
    \max_{k \in \{1,2\}}
    \left[
        \max(0, \cos(\move_{i,t}, \vect_{i,t}^{(k)}))
        \; s_{i,t}
        \; \sigma_{i,t}^{(k)}
    \right].
    \label{eq:frontier_reward}
\end{equation}
This term is deliberately secondary. It teaches directional preference toward
uncertain regions, while Eq.~\ref{eq:confidence_reward} remains the dominant
objective and determines whether the movement actually improves search.

\subsection{Survivor and Mission Rewards}

When UAV \(i\) scouts a previously unscouted survivor, it receives
\begin{equation}
    r^{\mathrm{scout}}_{i,t}
    =
    k_{\mathrm{scout}}
    \cdot n^{\mathrm{new\text{-}scout}}_{i,t}.
\end{equation}
All UAVs receive a smaller shared scout reward
\begin{equation}
    r^{\mathrm{team\text{-}scout}}_t
    =
    k_{\mathrm{team\text{-}scout}}
    \cdot n^{\mathrm{new\text{-}scout}}_{t}.
\end{equation}
Finally, when all five survivors have been scouted, all UAVs receive a one-time
mission-completion bonus \(r^{\mathrm{all\text{-}found}}_t\). These rewards use
actual detections, not survivor locations. We do not use progress shaping toward
survivors, because that would leak the unknown search target into the reward.

\subsection{Boundary Penalties}

We avoid action-level projection hacks that hide the reason for failure from the
policy. The simulator still enforces the physical domain for numerical validity,
but the learning signal is explicit. Let \(o_{i,t}\) be the fraction of UAV
\(i\)'s footprint outside the search area. We penalize outside footprint as
\begin{equation}
    c^{\mathrm{outside}}_{i,t}
    =
    k_{\mathrm{outside}} o_{i,t}^{2}.
    \label{eq:outside_penalty}
\end{equation}
Near the boundary, we also penalize actions that point outward:
\begin{equation}
    c^{\mathrm{outward}}_{i,t}
    =
    k_{\mathrm{outward}}
    o_{i,t}
    \max(0, \act_{i,t}^{\top}\mathbf{n}_{i,t})^2,
    \label{eq:outward_penalty}
\end{equation}
where \(\mathbf{n}_{i,t}\) is the outward normal of the nearest boundary. This
preserves freedom to move along the boundary while making repeated attempts to
leave the map consistently costly.

\section{Observations}

Each UAV receives a decentralized observation that includes ego state, teammate
state summaries, terrain context, and shared search memory. The search-memory
portion is expressed as uncertainty rather than coverage:
\begin{enumerate}[leftmargin=1.5em]
    \item \textbf{Global uncertainty map}: a pooled \(G \times G\) map of
    \(\unc_t(x)\), plus the mean confidence over the full map.
    \item \textbf{Local uncertainty patch}: a pooled \(L \times L\) patch around
    the UAV over a fixed physical radius, e.g. 150 m, computed from the same
    confidence grid.
    \item \textbf{Local max-uncertainty patch}: a companion max-pooled patch
    that preserves small high-uncertainty gaps that average pooling can wash out.
    \item \textbf{Top-2 frontier vectors}: the two sector-top frontier
    directions, distances, and scores from Eq.~\ref{eq:sector_score}.
    \item \textbf{Boundary state}: normalized distances to the four map edges
    and outside-footprint fraction.
    \item \textbf{Motion state}: current velocity, last displacement, previous
    action, footprint radius, and flight quality.
    \item \textbf{Team state}: relative UAV positions and shared frontier
    summaries in the communication-available setting.
\end{enumerate}
This observation design makes the reward locally interpretable: the policy sees
the same uncertainty field that generates the confidence reward and frontier
signal.

\section{Training Protocol}

We train in the UAV survivor diagnostic setting: three UAVs, zero UGVs, five
survivors, a 500 m by 500 m terrain, 300 simulation steps, and communication
available. Survivor and UAV start positions are sampled uniformly subject to UAV
spawn separation and edge-margin constraints, but survivor placement is not
conditioned on UAV placement. The actors share parameters within the UAV class,
which improves sample efficiency and prevents one drone from specializing into a
permanently bad role. The centralized critic observes the joint state and the
shared uncertainty memory.

The final system keeps the current empirically successful components--coverage
memory, top-2 frontier guidance, scout reward, shared UAV actors, and diagnostic
trajectory export--but changes the semantic unit of search from binary coverage
to probabilistic inspection confidence. This makes revisits neither always bad
nor always good: they are rewarded exactly when they increase the probability of
detecting a survivor.

\section{Diagnostics}

Evaluation uses deterministic and stochastic rollouts over fixed seed sets. We
report recall, final geometric coverage, final mean confidence, survivor scout
time, all-scouted time, path length, boundary exposure, same-step inter-UAV
overlap, and confidence-weighted revisit statistics. Survivor-level diagnostics
separate scouted and missed survivors by cumulative exposure probability, best
single-pass detection probability, normalized footprint distance, and central
versus edge exposure. Time-bin diagnostics report confidence gain, new binary
cells, frontier score, frontier reward, outside-footprint fraction, movement
speed, and reward/penalty scale over the episode.

The main success criterion is not coverage alone. A policy is considered
improved only if it increases recall or expected recall while maintaining high
confidence coverage and reducing avoidable high-confidence revisits. This
prevents reward hacking by broad but low-quality edge passes.

\section{Expected Advantages}

The confidence-aware formulation directly addresses the failure modes observed
in the UAV diagnostics:
\begin{enumerate}[leftmargin=1.5em]
    \item It aligns the dense reward with the probabilistic perception model,
    so a central high-probability pass is more valuable than an edge skim.
    \item It allows useful revisits of low-confidence cells, avoiding the
    brittle behavior produced by purely binary overlap penalties.
    \item It assigns marginal credit when UAV footprints overlap, reducing the
    need for hand-tuned inter-UAV overlap penalties.
    \item It makes late-episode cleanup meaningful, because sparse remaining
    uncertainty can still produce confidence reward.
    \item It preserves the successful frontier idea but computes frontiers over
    uncertainty mass instead of binary uncovered area.
\end{enumerate}

\section{Limitations and Communication Dropout}

This system assumes a communication-available scenario in which all UAVs update
and observe the same confidence map. Under communication dropout, each UAV would
continue to update a private confidence map and the team confidence update in
Eq.~\ref{eq:confidence_update} would no longer be globally available. The same
reward structure can be used locally, but marginal credit and frontier ownership
must be computed from each UAV's belief state until communication is restored.
We therefore treat dropout as a separate extension rather than mixing it into
the communication-available diagnostic baseline.

\end{document}
